%%% SECTION 2: U.S. GOVERNMENT FRAMEWORK FOR PHYSICAL AI TRIALS %%%

Physical AI oncology trials must be understood within the full constitutional
and statutory framework of U.S. government, spanning the Legislative, Executive,
and Judicial branches. The governance structure for Physical AI builds on
existing federal branch operations rather than requiring new institutional
frameworks, leveraging established institutional credibility and legal authority.
This section adapts the comprehensive research on U.S. federal branch operations
statutes and their application to AI-enabled clinical trials to Physical AI
oncology trial contexts.

\subsection{Constitutional Foundation}
\label{subsec:constitutional}

The U.S. federal government's three-branch structure is established primarily by
the U.S. Constitution. Article I vests legislative powers in a Congress
consisting of a Senate and House of Representatives. Article II vests executive
power in a President, with the Appointments Clause establishing Senate
advice-and-consent for officers and the Take Care duty (Article II, Section 3)
providing the constitutional anchor for executive implementation and supervision
of federal law. Article III vests the judicial power in one Supreme Court and
such inferior courts as Congress may establish.

For Physical AI oncology trials, this three-branch structure provides essential
checks and balances. Legislative authority creates the statutory framework
through which Physical AI trial programs are funded and overseen. Executive
authority implements that framework through FDA and HHS, which directly
regulate clinical trial conduct. Judicial authority resolves disputes and
interprets the boundaries of regulatory action. The Congressional Review Act
(5 U.S.C. \S~801 et seq.) provides legislative oversight of executive
rulemaking, meaning that any FDA regulations specific to Physical AI would be
subject to congressional review. The Administrative Procedure Act (5 U.S.C.
\S~551 et seq.) requires that Physical AI regulations follow standard
notice-and-comment rulemaking procedures, ensuring public participation and
transparency.

The Appointments Clause has practical implications for Physical AI governance:
FDA commissioners and senior officials who would oversee Physical AI trial
approvals are appointed through the constitutionally mandated process, ensuring
democratic accountability for regulatory decisions affecting patient safety.
The Take Care duty requires the executive branch to faithfully implement
clinical trial statutes, preventing either over-regulation that would block
beneficial Physical AI innovations or under-regulation that would compromise
patient safety.

\subsection{Legislative Branch and Physical AI Regulation}
\label{subsec:legislative}

The legislative branch's role in Physical AI regulation operates through several
operational statutes. The Congressional Budget and Impoundment Control Act of
1974 (codified across titles, with CBO at 2 U.S.C. \S~601) creates the
Congressional Budget Office and statutory budget procedures that structure
annual budgeting and congressional scorekeeping. FDA funding for Physical AI
trial review capacity depends directly on congressional appropriations authority
exercised through these budget procedures.

The Congressional Review Act (5 U.S.C. \S~801 et seq.) establishes a mechanism
for Congress to review and disapprove certain agency rules under defined
procedures. If FDA were to promulgate rules specifically governing Physical AI
clinical trials, these rules would be subject to congressional review,
providing a check on executive rulemaking power. The Congressional
Accountability Act of 1995 (2 U.S.C. ch. 24, \S~1301; Pub. L. 104-1) applies
workplace and labor protections to legislative-branch employing offices,
establishing precedent for how new categories of workers (including those
overseeing Physical AI systems) receive legislative protection.

The Government Accountability Office (GAO), with Comptroller General
authorities established under 31 U.S.C. \S~702, provides audit and
investigation capacity for oversight of Physical AI trial programs. GAO audits
of FDA's review processes for AI-enabled medical products would apply equally
to Physical AI trial oversight, ensuring that congressional appropriations are
used effectively and that regulatory processes meet statutory requirements.

Congressional appropriations directly affect FDA's capacity to review and
approve Physical AI trial applications. If Physical AI trials require new
categories of expertise (robotics engineers, AI safety specialists,
cybersecurity analysts) within FDA review teams, congressional funding must
support these positions. The budget process statutes ensure that these resource
needs are evaluated through established scorekeeping procedures rather than
ad hoc allocation.

\begin{table}[H]
\centering
\caption{Core Legislative Statutes Applicable to Physical AI Trial Governance}
\label{tab:legislative-statutes}
\small
\begin{tabularx}{\textwidth}{l X l}
\toprule
\textbf{Statute} & \textbf{Physical AI Relevance} & \textbf{Citation} \\
\midrule
Budget and Impoundment Control Act & FDA funding for Physical AI review capacity & 2 U.S.C. \S~601 \\
Congressional Review Act & Review of Physical AI-specific FDA rules & 5 U.S.C. \S~801 \\
Congressional Accountability Act & Worker protections for Physical AI oversight staff & 2 U.S.C. \S~1301 \\
GAO/Comptroller General & Audit of Physical AI trial programs & 31 U.S.C. \S~702 \\
Reapportionment Act & House composition for Physical AI policy votes & 2 U.S.C. \S~2a \\
\bottomrule
\end{tabularx}
\end{table}

\subsection{Executive Branch Implementation}
\label{subsec:executive}

The executive branch implements clinical trial regulation through FDA and HHS,
with several foundational statutes governing the process. The Administrative
Procedure Act (APA, Pub. L. 79-404; 5 U.S.C. \S~551 et seq.) defines
``agency'' and related terms, forming the baseline procedural code for federal
administrative action and review. FDA's authority to regulate Physical AI trials
derives from the APA's rulemaking framework: any Physical AI-specific
regulations would follow notice-and-comment procedures, with judicial review
available for arbitrary or capricious agency action.

The Freedom of Information Act (FOIA, Pub. L. 89-487; 5 U.S.C. \S~552)
creates enforceable disclosure duties subject to enumerated exemptions. For
Physical AI trials, FOIA ensures that FDA review processes, approval decisions,
and safety data are accessible to the public, researchers, and industry
participants. The transparency provided by FOIA supports public trust in
Physical AI adoption by making regulatory decision-making visible and
accountable.

The Privacy Act of 1974 (Pub. L. 93-579; 5 U.S.C. \S~552a) establishes
restrictions on disclosure and requirements for systems of records maintained by
federal agencies. When FDA maintains records about Physical AI trial sponsors,
investigators, or participants, the Privacy Act governs how those records may
be used and disclosed. This statute includes criminal penalties for certain
willful violations, providing strong incentives for proper data handling.

The Federal Register Act (44 U.S.C. \S~1505) requires publication of
presidential proclamations and executive orders with general applicability and
legal effect. Any executive orders directing FDA to prioritize Physical AI
trial review or establishing Physical AI governance frameworks must be
published in the Federal Register to have legal effect. The Paperwork Reduction
Act (44 U.S.C. \S~3501) requires OMB review of information collections,
meaning that new data collection requirements for Physical AI trial sponsors
must go through OMB clearance.

The Inspectors General framework (5 U.S.C. ch. 4) requires IG reporting
structures within executive agencies, providing internal oversight and audit
capability for Physical AI trial programs. The Federal Vacancies Reform Act of
1998 (5 U.S.C. \S~3345) establishes rules for acting officials, ensuring
continuity of FDA leadership during transitions that might affect Physical AI
regulatory decisions.

\begin{table}[H]
\centering
\caption{Core Executive Statutes Applicable to Physical AI Trial Operations}
\label{tab:executive-statutes}
\small
\begin{tabularx}{\textwidth}{l X l}
\toprule
\textbf{Statute} & \textbf{Physical AI Relevance} & \textbf{Citation} \\
\midrule
Administrative Procedure Act & Rulemaking for Physical AI regulations & 5 U.S.C. \S~551 \\
Freedom of Information Act & Transparency of Physical AI trial decisions & 5 U.S.C. \S~552 \\
Privacy Act of 1974 & Protection of Physical AI trial records & 5 U.S.C. \S~552a \\
Federal Register Act & Publication of Physical AI executive orders & 44 U.S.C. \S~1505 \\
Paperwork Reduction Act & OMB review of Physical AI data collection & 44 U.S.C. \S~3501 \\
Inspectors General & Internal audit of Physical AI programs & 5 U.S.C. ch. 4 \\
Federal Vacancies Reform & Continuity of FDA Physical AI leadership & 5 U.S.C. \S~3345 \\
Presidential Records Act & Preservation of Physical AI policy records & 44 U.S.C. \S~2201 \\
\bottomrule
\end{tabularx}
\end{table}

The three adapted regulatory standards in this platform demonstrate how
existing executive branch regulatory frameworks can be adapted rather than
replaced. The Adaption: ICH Harmonised Guideline \cite{ich-adapt} builds on the
internationally recognized GCP standard. The Physical AI Adaption: 21 CFR Part
50 \cite{cfr50-adapt} extends existing human subjects protection regulation. The
Physical AI Adaption: 21 CFR Part 312 \cite{cfr312-adapt} extends the IND
application regulation. Each adaptation operates within the APA's rulemaking
framework, demonstrating that Physical AI can be regulated through existing
administrative procedures.

\subsection{Judicial Branch Oversight}
\label{subsec:judicial}

The judicial branch provides the framework for resolving Physical AI regulatory
disputes. The modern statutory backbone for court organization, jurisdiction,
and judicial administration is established in Title 28 of the U.S. Code. The
Supreme Court size and quorum are set by 28 U.S.C. \S~1, with circuit
composition defined by 28 U.S.C. \S~41 (currently thirteen circuits including
the Federal Circuit). The Judicial Conference (28 U.S.C. \S~331) serves as the
judiciary's policy-making body, coordinating administration across all federal
courts.

The Rules Enabling Act (28 U.S.C. \S~2072) authorizes the Supreme Court to
prescribe general rules of practice, procedure, and evidence, while prohibiting
rules that abridge, enlarge, or modify substantive rights. These procedural
rules govern how Physical AI regulatory disputes would be litigated, from
initial filing through appeal. The judicial discipline and complaints process
(28 U.S.C. \S~351) ensures accountability for judges who handle Physical AI
cases.

Courts would review FDA decisions regarding Physical AI trial approvals, safety
actions, and clinical holds under existing administrative law frameworks. The
APA provides for judicial review of agency action, with courts evaluating
whether FDA's Physical AI decisions are arbitrary, capricious, an abuse of
discretion, or otherwise not in accordance with law. This standard of review
means that FDA must articulate reasoned explanations for Physical AI regulatory
decisions, which in turn requires the kind of comprehensive evidence base that
this National Platform provides.

If FDA were to place a clinical hold on a Physical AI trial, the sponsor could
challenge that decision through administrative channels and ultimately through
judicial review. If FDA were to deny an IND application that included Physical
AI components, the denial would be subject to review under the same standards
that apply to all IND decisions. The existing judicial framework thus provides
a safety valve against both over-regulation and under-regulation of Physical AI
trials, ensuring that regulatory decisions are grounded in evidence rather than
speculation.

\subsection{Cross-Branch Governance for Physical AI}
\label{subsec:cross-branch}

Several statutes operate across all three branches to ensure accountability in
Physical AI trial governance. Transparency and recordkeeping requirements
include Federal Register publication requirements (44 U.S.C. \S~1505), FOIA
public access provisions (5 U.S.C. \S~552), agency record preservation
obligations (44 U.S.C. \S~3101), and presidential record definitions
(44 U.S.C. \S~2201). Together, these ensure that Physical AI policy decisions
are documented, preserved, and accessible for public scrutiny.

Budget execution constraints apply through the Antideficiency Act (31 U.S.C.
\S~1341), which prohibits obligations and expenditures exceeding available
appropriations. This means that Physical AI trial programs cannot spend beyond
congressional authorization, providing fiscal discipline for potentially
expensive national deployment efforts. The Ethics in Government Act framework
establishes disclosure and ethics requirements spanning all branches, ensuring
that officials involved in Physical AI regulatory decisions are subject to
conflict-of-interest protections and financial disclosure requirements.

Physical AI oncology trials benefit from the full weight of these existing
cross-branch governance mechanisms. Rather than creating new governance
structures, this platform adapts proven mechanisms to ensure transparency,
fiscal responsibility, and ethical conduct. This approach is consistent with the
platform's broader philosophy of building on established frameworks rather than
replacing them.

\subsection{Current Regulatory Framework for AI in Clinical Trials}
\label{subsec:current-ai-framework}

AI use in clinical trials is regulated primarily through what the AI does rather
than through a dedicated ``AI clinical trials'' statute. FDA guidance
distinguishes four categories of AI use in trial contexts. First, AI used to
generate information or data intended to support regulatory decision-making for
drugs and biologics falls under FDA's January 2025 draft guidance, which
proposes a risk-based credibility assessment framework with seven steps: define
question of interest, define context of use, assess model risk, plan, execute,
document, and determine adequacy. Second, AI embedded in or functioning as
medical device software falls under device lifecycle management guidance and
PCCP requirements. Third, AI used as part of trial conduct and data systems
must comply with 21 CFR Part 11 electronic records requirements and related
computerized systems guidance. Fourth, AI as human subjects research and data
risk falls under Common Rule and IRB review requirements.

\begin{table}[H]
\centering
\caption{Federal Regulations and Guidances Relevant to Physical AI Trials}
\label{tab:regulations-ai}
\small
\begin{tabularx}{\textwidth}{l l X}
\toprule
\textbf{Type} & \textbf{Citation} & \textbf{Physical AI Relevance} \\
\midrule
Regulation & 21 CFR Part 312 & IND for trials with Physical AI drug administration \\
Regulation & 21 CFR Part 50 & Informed consent for Physical AI robotic procedures \\
Regulation & 21 CFR Part 56 & IRB review of Physical AI system risks \\
Regulation & 21 CFR Part 812 & IDE for Physical AI device investigations \\
Regulation & 21 CFR Part 11 & Electronic records from Physical AI systems \\
Regulation & 42 CFR Part 11 & Trial registration including Physical AI methods \\
Regulation & 45 CFR Part 46 & Common Rule for Physical AI human subjects research \\
Regulation & 45 CFR Parts 160, 164 & HIPAA for Physical AI health data \\
Guidance & AI Drug/Biologic (Jan 2025) & AI credibility framework for Physical AI evidence \\
Guidance & AI Device Lifecycle (Jan 2025) & Lifecycle management for Physical AI software \\
Guidance & Decentralized Trials (Sep 2024) & Remote Physical AI monitoring elements \\
Guidance & PCCP for AI Devices (Aug 2025) & Change control for Physical AI system updates \\
Guidance & CDS Software (Jan 2026) & Classification of Physical AI decision support \\
\bottomrule
\end{tabularx}
\end{table}

Physical AI introduces additional considerations across all four categories.
When Physical AI systems generate evidence for regulatory decisions (such as
robotic surgical outcomes data), the credibility assessment must account for
robot-specific variables including mechanical precision, sensor accuracy, and
autonomous decision quality. When Physical AI functions as a regulated device,
the lifecycle management must address hardware degradation, calibration drift,
and physical wear in addition to software updates. When Physical AI generates
electronic trial records (robot telemetry, sensor data, autonomous decision
logs), these records represent new categories that must comply with Part 11
requirements. When Physical AI interacts directly with human subjects, the risk
profile differs fundamentally from software-only AI because of physical contact
and proximity.

\subsection{AI and Oncology: A High-Frequency Use Case}
\label{subsec:ai-oncology}

The FDA Oncology Center of Excellence \cite{fda-oncology-guidance} describes an
Oncology AI Program begun in 2023, linking AI-related FDA guidances used in
oncology regulatory contexts. In oncology trials specifically, AI is frequently
implicated because endpoints and eligibility often rely on imaging (CT, MRI,
PET), digital pathology, complex biomarker patterns, and longitudinal
EHR-derived outcomes, all of which are common AI input modalities.

Physical AI extends these existing AI applications by adding robotic
embodiment. Rather than AI analyzing images to inform human decisions, Physical
AI systems can directly act on imaging analysis by guiding surgical robots to
tumor sites, positioning radiotherapy beams with sub-millimeter precision, or
directing needle-placement robots to biopsy targets. A Cancer Patient's
Journey, Physical AI \cite{patient-journey-paper} documents how this integrated
approach works across all ten stages of an oncology trial, from prescreening
through closeout. The ten-stage pipeline demonstrates that Physical AI is not
limited to isolated procedures but can manage the complete patient experience.

\subsection{Practical Compliance Synthesis}
\label{subsec:compliance-synthesis}

In oncology trials using Physical AI, compliance requires building a defensible
chain from context of use through validation, auditability, and human subjects
protection.

If Physical AI outputs will be used as evidence for FDA decisions about safety,
effectiveness, or quality, sponsors should document the model's context of use,
risk, and credibility assessment consistent with FDA's draft framework. The
adapted 21 CFR Part 312 \cite{cfr312-adapt} provides the IND framework for
including Physical AI evidence in regulatory submissions, while the adapted ICH
E6(R3) \cite{ich-adapt} establishes the quality management system for ensuring
Physical AI data integrity.

If Physical AI functions as a regulated device, the trial may implicate IDE and
device marketing pathways, with lifecycle management and PCCP expectations
governing what changes are allowed and how they are controlled. The USL scoring
framework \cite{usl-standard} provides quantitative assessment of each robot's
readiness for clinical deployment, with Dimension D (clinical trial
collaboration) directly measuring regulatory compliance status.

If Physical AI relies on decentralized or digital health technology data, the
data attribution, metadata integrity, and Part 11 boundaries become critical
compliance points. The MCP server infrastructure \cite{national-mcp-paper}
addresses this through standardized data flow pathways and conformance level
requirements that ensure data integrity across distributed Physical AI systems.

If Physical AI training and validation rely on secondary data and research
exemptions, IRB considerations around identifiability and bias must be
addressed. The federated learning framework \cite{fl-paper} provides
privacy-preserving approaches to model training that reduce identifiability
risks while enabling multi-site data utilization under HIPAA \cite{hipaa} and
state privacy law requirements.

The adapted 21 CFR Part 50 \cite{cfr50-adapt} ensures that human subjects
protection extends to all Physical AI interactions, with informed consent
requirements covering robot descriptions, AI decision-making transparency,
patient rights to refuse robotic procedures, and emergency stop procedures.
Together, the three adapted standards provide comprehensive responses to each
compliance pathway, demonstrating that Physical AI can operate within the
existing regulatory architecture with targeted adaptations.

%%% ADDITIONAL CONTENT FOR SECTION 2 %%%

\subsubsection{Detailed FDA Guidance Applicability to Physical AI}

The current FDA guidance landscape relevant to Physical AI clinical trials
encompasses more than a dozen individual guidance documents, each addressing
a distinct regulatory concern. No single guidance document was written with
Physical AI in mind, yet each contains provisions that apply to one or more
aspects of Physical AI deployment in oncology trials. Table~\ref{tab:fda-guidance-pai}
provides a detailed mapping of major FDA guidance documents against six
Physical AI applicability dimensions: robotic hardware, AI software, data
systems, patient interaction, trial operations, and multi-site coordination.

\clearpage

\begin{table}[H]
\centering
\caption{FDA Guidance Documents Mapped to Physical AI Applicability Dimensions}
\label{tab:fda-guidance-pai}
\small
\begin{tabularx}{\textwidth}{X c c c c c c}
\toprule
\textbf{FDA Guidance} & \rotatebox{70}{\textbf{Robotic HW}} & \rotatebox{70}{\textbf{AI Software}} & \rotatebox{70}{\textbf{Data Systems}} & \rotatebox{70}{\textbf{Patient Int.}} & \rotatebox{70}{\textbf{Trial Ops}} & \rotatebox{70}{\textbf{Multi-Site}} \\
\midrule
AI for Drug/Biologic Decisions (Jan 2025) & - & Yes & Yes & - & Yes & - \\
AI Device Lifecycle (Jan 2025) & Yes & Yes & - & - & - & - \\
Decentralized Trials (Sep 2024) & - & - & Yes & Yes & Yes & Yes \\
PCCP for AI Devices (Aug 2025) & Yes & Yes & - & - & - & - \\
CDS Software (Jan 2026) & - & Yes & - & Yes & - & - \\
Adaptive Designs (2019) & - & - & - & - & Yes & Yes \\
Electronic Records (Part 11) & - & - & Yes & - & Yes & Yes \\
RAS Safety (2015) & Yes & - & - & Yes & - & - \\
Cybersecurity Medical Devices (2023) & Yes & Yes & Yes & - & - & Yes \\
Patient Reported Outcomes (2009) & - & - & Yes & Yes & Yes & - \\
Risk-Based Monitoring (2013) & - & - & Yes & - & Yes & Yes \\
Computer System Validation (2020) & - & Yes & Yes & - & Yes & - \\
\bottomrule
\end{tabularx}
\end{table}

The mapping in Table~\ref{tab:fda-guidance-pai} reveals a structural gap in
the current guidance landscape. No single guidance document addresses more
than three of the six Physical AI applicability dimensions simultaneously.
The most broadly applicable guidance, the decentralized trials guidance
(September 2024), covers four dimensions but excludes robotic hardware and
AI software, which are core to Physical AI. The AI device lifecycle
guidance (January 2025) addresses robotic hardware and AI software but
excludes data systems, patient interaction, trial operations, and multi-site
coordination. This fragmentation means that a Physical AI trial sponsor must
synthesize requirements from at least six to eight separate guidance
documents to achieve comprehensive compliance.

The adapted regulatory standards in this National Platform
\cite{ich-adapt, cfr50-adapt, cfr312-adapt} address all six dimensions
simultaneously, providing the unified reference that the current guidance
landscape lacks. The adapted ICH E6(R3) \cite{ich-adapt} addresses trial
operations, data systems, and patient interaction through its quality
management, data governance, and informed consent provisions. The adapted
21 CFR Part 312 \cite{cfr312-adapt} addresses robotic hardware, AI software,
and multi-site coordination through its Physical AI system descriptions,
validation requirements, and multi-site IND provisions. The adapted 21 CFR
Part 50 \cite{cfr50-adapt} addresses patient interaction with comprehensive
informed consent requirements that cover all ten Physical AI system
categories.

The AI for Drug/Biologic Regulatory Decisions guidance (January 2025) merits
particular attention because its seven-step credibility assessment framework
provides the closest existing analog to the validation approach required for
Physical AI evidence. The seven steps are: (1) define the question of
interest, (2) define the context of use, (3) assess model risk, (4) develop
a credibility assessment plan, (5) execute the plan, (6) document the
assessment, and (7) determine adequacy. For Physical AI systems, each step
requires expansion to account for hardware-software integration, real-time
adaptive behavior, and multi-system coordination. Step 3, model risk
assessment, is particularly complex for Physical AI because risk encompasses
not only AI model prediction errors but also robot mechanical failure, sensor
degradation, communication latency, and cybersecurity compromise. The
adapted 21 CFR Part 312 \cite{cfr312-adapt} provides this expanded risk
framework through its Phase 0 simulation validation requirement, which
mandates comprehensive system-level risk assessment before human subjects
enrollment.

The Predetermined Change Control Plan (PCCP) guidance (August 2025)
addresses how AI-enabled devices can be modified after initial marketing
authorization without requiring new regulatory submissions for each change.
For Physical AI in clinical trials, the PCCP concept is directly relevant
because AI algorithms improve through learning during the trial itself. The
federated learning framework \cite{fl-paper} implements a PCCP-consistent
approach through its triple AI peer review pipeline, where model updates are
validated by three independent AI systems (Claude Opus 4.6 \cite{claude-opus},
GPT-5.4 \cite{gpt-5-4}, and Gemini \cite{gemini}) before deployment across
trial sites. This triple review process provides the ``pre-specified
modification protocol'' that the PCCP guidance envisions, adapted for the
multi-site, continuous-learning context of Physical AI trials.

The decentralized trials guidance (September 2024) addresses remote
monitoring, digital health technology integration, and decentralized trial
elements that are directly applicable to Physical AI multi-site operations.
Physical AI extends the decentralized trial concept by enabling fully
autonomous site operations with remote physician oversight, continuous remote
data collection through robot-mounted sensors, and real-time inter-site
coordination through MCP servers \cite{national-mcp-paper}. However, the
decentralized trials guidance was written primarily for software-based
digital health technologies and does not address the physical presence of
robotic systems at trial sites, creating a gap that the site establishment
documentation \cite{site-docs} fills with its eleven regulatory documents
covering building standards, operational procedures, and emergency
preparedness.

The cybersecurity guidance for medical devices (2023) addresses
authentication, encryption, software bill of materials, and vulnerability
management for connected medical devices. Physical AI systems are inherently
connected devices that communicate with central servers, other robots, and
clinical data systems through network protocols. The cybersecurity
requirements established in the MCP server architecture
\cite{national-mcp-paper} exceed the minimum standards in the FDA
cybersecurity guidance by implementing five conformance levels with
progressive security requirements, mandatory encryption for all data
transmissions, and real-time intrusion detection through the MCP safety
server.

\subsubsection{Guidance Gap Analysis and Platform Response}

The comprehensive mapping of FDA guidances to Physical AI dimensions reveals
three categories of gaps. First, coverage gaps exist where no current
guidance addresses a specific Physical AI concern. The most significant
coverage gap is the absence of guidance for multi-robot coordination in
clinical settings, where multiple Physical AI systems must share workspace,
exchange data, and coordinate patient care activities. The MCP server
architecture \cite{national-mcp-paper} addresses this gap directly.

Second, depth gaps exist where current guidance addresses a Physical AI
concern but with insufficient specificity. The AI device lifecycle guidance
addresses software updates but does not specify validation requirements for
simultaneous hardware and software updates, which are common in Physical AI
systems where firmware upgrades may change robot motion profiles. The adapted
ICH E6(R3) \cite{ich-adapt} addresses this through its system documentation
requirements in Appendix A, which mandate validation of all system changes
including hardware-software interactions.

Third, integration gaps exist where multiple guidances each address part of
a Physical AI concern but no guidance addresses the integrated whole. The
intersection of informed consent (addressed by multiple guidances), device
safety (addressed by RAS and cybersecurity guidances), and data integrity
(addressed by Part 11 and risk-based monitoring guidances) creates a complex
compliance landscape that no single guidance navigates. The three adapted
standards in this platform \cite{ich-adapt, cfr50-adapt, cfr312-adapt}
resolve this integration gap by providing cross-referenced, internally
consistent requirements that address the integrated Physical AI system
rather than its individual components.

\subsection{Federal Agency Coordination for Physical AI}
\label{subsec:federal-agency-coordination}

Physical AI oncology trials require coordinated oversight from multiple
federal agencies, each with distinct statutory authorities, regulatory
mandates, and institutional expertise. Effective coordination among these
agencies is essential for ensuring that Physical AI trials proceed without
duplicative requirements, conflicting guidance, or regulatory gaps that
could compromise patient safety or impede innovation. This subsection
analyzes the coordination requirements among the four primary federal
agencies with jurisdiction over Physical AI oncology trials: the Food and
Drug Administration (FDA), the Department of Health and Human Services (HHS),
the National Institutes of Health (NIH), and the National Cancer Institute
(NCI).

\subsubsection{FDA Organizational Structure for Physical AI}

FDA's regulatory authority over Physical AI trials spans multiple centers
within the agency. The Center for Drug Evaluation and Research (CDER)
reviews IND applications for investigational drugs administered by Physical
AI systems. The Center for Devices and Radiological Health (CDRH) reviews
the Physical AI robotic devices themselves, including surgical robots,
collaborative robots, and radiotherapy positioning systems. The Center for
Biologics Evaluation and Research (CBER) reviews biologics administered
through Physical AI systems. The Oncology Center of Excellence (OCE)
\cite{fda-oncology-guidance} provides oncology-specific expertise and
coordinates cross-center review of oncology products.

The multi-center nature of Physical AI review creates coordination
challenges within FDA itself. A Physical AI oncology trial involving a
robotic system (CDRH jurisdiction) administering an investigational drug
(CDER jurisdiction) for a biologic-sensitive cancer type (CBER interest) in
an oncology context (OCE coordination) requires input from at least four
organizational units. FDA's existing intercenter agreement and lead center
designation processes address this through assignment of a lead review
center, but these processes were not designed for the level of cross-center
integration that Physical AI requires. The adapted 21 CFR Part 312
\cite{cfr312-adapt} facilitates this coordination by providing a unified
IND framework that addresses device, drug, and operational requirements in a
single submission, reducing the need for separate cross-center negotiations.

\subsubsection{HHS Oversight and Policy Coordination}

The Department of Health and Human Services provides overarching policy
direction for Physical AI governance through several mechanisms. The Office
for Human Research Protections (OHRP) within HHS administers the Common Rule
(45 CFR Part 46), which establishes baseline human subjects protections for
all federally funded research. Physical AI trials funded through NIH grants
or conducted at institutions receiving HHS funding must comply with Common
Rule requirements in addition to FDA-specific protections under 21 CFR Part
50 \cite{cfr-part-50}.

The Office of the National Coordinator for Health Information Technology
(ONC) within HHS sets standards for health information exchange that affect
how Physical AI systems share data with electronic health records, laboratory
information systems, and clinical trial management systems. ONC's adoption of
FHIR \cite{hl7-fhir} as the standard for health data interoperability
provides the foundation for Physical AI data exchange, extended by the MCP
server infrastructure \cite{national-mcp-paper} to accommodate Physical
AI-specific data types including robot telemetry and autonomous decision logs.

The Office for Civil Rights (OCR) within HHS enforces HIPAA \cite{hipaa}
privacy and security requirements that apply to all Physical AI systems
handling protected health information. OCR's enforcement actions and guidance
documents define the boundaries of permissible data use, breach notification
requirements, and security standards that Physical AI systems must meet. The
federated learning framework \cite{fl-paper} was designed with HIPAA
compliance as a foundational requirement, implementing privacy-preserving
approaches that minimize the exposure of protected health information during
multi-site model training.

HHS also coordinates cross-agency policy development through initiatives
such as the HHS AI Strategy and the National AI Research Resource. These
policy frameworks establish priorities for AI research and deployment that
affect Physical AI trial development. The Physical AI National Platform
aligns with HHS policy priorities by demonstrating how AI can improve health
outcomes (through more efficient clinical trials), reduce costs (through
operational efficiencies documented in the financial analysis), and protect
patient rights (through the adapted informed consent framework).

\subsubsection{NIH Research Coordination}

The National Institutes of Health plays a dual role in Physical AI oncology
trials: as a funder of research that may utilize Physical AI systems and as
a standard-setter for clinical trial conduct. NIH clinical trial policies
require registration on ClinicalTrials.gov \cite{clinicaltrials-gov}, results
reporting, and adherence to good clinical practice standards. Physical AI
trials funded by NIH must comply with these policies while also addressing
the unique requirements of robotic systems in clinical settings.

NIH's clinical trial stewardship framework, which defines expectations for
NIH-funded clinical trial networks, has implications for Physical AI trial
design and oversight. The framework emphasizes adaptive trial designs
\cite{fda-adaptive-designs}, centralized monitoring, and data sharing, all
of which are enhanced by Physical AI capabilities. The MCP server
architecture \cite{national-mcp-paper} provides the technical infrastructure
for the centralized monitoring and data sharing that NIH's stewardship
framework envisions, while the federated learning framework \cite{fl-paper}
enables the multi-site data analysis that adaptive designs require.

NIH's Artificial Intelligence and Technology Collaboratories (AITC) program
and related initiatives provide research infrastructure for evaluating AI
technologies in healthcare settings. Physical AI trial sponsors may leverage
these NIH resources for system validation, training data development, and
clinical evidence generation. The open-source nature of the National Platform
repositories \cite{main-repo, mcp-repo, fl-repo} facilitates integration
with NIH research infrastructure by providing transparent, auditable code
that NIH-funded investigators can evaluate and extend.

\subsubsection{NCI Trial Network Integration}

The National Cancer Institute operates the National Clinical Trials Network
(NCTN), the Community Oncology Research Program (NCORP), and the NCI Central
Institutional Review Board (CIRB) \cite{nci-cirb}, which together constitute
the primary infrastructure for federally funded oncology clinical trials in
the United States. Physical AI integration with these NCI networks requires
coordination across multiple organizational boundaries.

The NCI CIRB provides centralized IRB review for NCI-sponsored multi-site
trials, reducing the administrative burden of obtaining separate IRB approval
at each participating site. For Physical AI trials, centralized CIRB review
offers a significant advantage: rather than requiring each site's IRB to
independently evaluate Physical AI system risks (which may require robotics
and AI expertise not available at every institution), the CIRB can develop
specialized Physical AI review capability and apply it consistently across
all participating sites. The adapted 21 CFR Part 50 \cite{cfr50-adapt}
provides the informed consent framework that the CIRB would use to evaluate
Physical AI-specific consent requirements.

NCI's modernization efforts \cite{nci-modernizing} explicitly acknowledge the
need for technological innovation in clinical trial infrastructure. Physical
AI aligns with NCI's modernization priorities by addressing enrollment
challenges (through expanded operational hours and AI-assisted screening),
data quality concerns (through automated electronic capture), and multi-site
coordination difficulties (through MCP servers and federated learning). The
National Platform's phased implementation strategy
(Section~\ref{sec:implementation}) can be aligned with NCI's modernization
timeline to ensure complementary rather than competing infrastructure
development.

\subsubsection{Interagency Coordination Mechanisms}

Effective Physical AI governance requires coordination mechanisms that bridge
the organizational boundaries among FDA, HHS, NIH, and NCI. Several existing
mechanisms provide frameworks for this coordination.

\begin{table}[H]
\centering
\caption{Federal Agency Coordination for Physical AI Trial Oversight}
\label{tab:agency-coordination}
\small
\begin{tabularx}{\textwidth}{l l X}
\toprule
\textbf{Agency} & \textbf{Primary Authority} & \textbf{Physical AI Coordination Role} \\
\midrule
FDA (CDER) & IND review (21 CFR 312) & Review Physical AI drug administration protocols \\
FDA (CDRH) & Device review (21 CFR 812) & Evaluate Physical AI robotic device safety \\
FDA (OCE) & Oncology coordination & Cross-center review of Physical AI oncology applications \\
HHS (OHRP) & Common Rule (45 CFR 46) & Human subjects oversight for Physical AI research \\
HHS (ONC) & Health IT standards & Interoperability standards for Physical AI data \\
HHS (OCR) & HIPAA enforcement & Privacy and security for Physical AI health data \\
NIH & Research funding/policy & Fund Physical AI research and set trial standards \\
NCI & Cancer trial networks & Integrate Physical AI into NCTN/NCORP infrastructure \\
NCI CIRB & Centralized IRB review & Specialized Physical AI protocol review \\
\bottomrule
\end{tabularx}
\end{table}

The Reagan-Udall Foundation, established by Congress as a non-profit
foundation supporting FDA's mission, provides a mechanism for convening
stakeholders across government, industry, and academia to address emerging
regulatory challenges. A Reagan-Udall-convened working group on Physical AI
clinical trials could bring together representatives from all relevant
agencies to develop coordinated guidance, share expertise, and identify
regulatory gaps.

Interagency memoranda of understanding (MOUs) provide formal coordination
frameworks between specific agencies. An MOU between FDA and NCI addressing
Physical AI trial oversight could establish shared review procedures, mutual
recognition of Physical AI system evaluations, and coordinated safety
monitoring responsibilities. The existing FDA-NCI collaboration on the
Oncology Center of Excellence \cite{fda-oncology-guidance} provides a model
for such coordination.

Congressional mandates through the FD\&C Act reauthorization process (most
recently the FDA Reauthorization Act) periodically update FDA's authorities
and mandates, providing opportunities to incorporate Physical AI-specific
provisions. The Congressional Budget and Impoundment Control Act, referenced
in Table~\ref{tab:legislative-statutes}, ensures that any new Physical AI
oversight responsibilities are accompanied by appropriate funding
authorization.

The Physical AI National Platform facilitates interagency coordination by
providing a common reference framework that all agencies can use. Rather than
each agency developing independent Physical AI evaluation criteria, the
adapted standards \cite{ich-adapt, cfr50-adapt, cfr312-adapt}, the
quantitative scoring frameworks \cite{usl-standard}, and the infrastructure
specifications \cite{national-mcp-paper, fl-paper} provide a shared
vocabulary and shared evidence base that reduce coordination costs and
prevent inconsistent requirements. This unifying function may be the
platform's most significant contribution to the governmental framework for
Physical AI trials.

%%% END ADDITIONAL CONTENT FOR SECTION 2 %%%
